% ============================================================
%  Chapter 11: Relativity
% ============================================================
\chapter{相对论：从光速不变到引力}
\label{ch:relativity}

\section{概述}

R1（因果时空结构）的直接推论涵盖狭义相对论的运动学（Lorentz 变换、质能等价）和广义相对论（Einstein 场方程）。本章从 R1 的子组件出发，严格推导这些结果。

% ════════════════════════════════════════════════════════
\section{Lorentz 变换}
% ════════════════════════════════════════════════════════

\begin{lawbox}[Lorentz 变换]
  $t' = \gamma(t - vx/c^2)$,\quad $x' = \gamma(x - vt)$,\quad $\gamma = 1/\sqrt{1 - v^2/c^2}$
  \quad\textbf{前提}：c 不变 (R1), Lorentz 不变性 (R1)
\end{lawbox}

\noindent\textbf{变量}：$\gamma \ge 1$ [1]（Lorentz 因子），$v$ [L$\cdot$T$^{-1}$]（相对速度，$|v| < c$）。

\begin{derivationbox}[推导]
  \textbf{Step 1} — 基本假设：$c' = c$（R1: kernel.light\_speed\_invariance）

  \textbf{Step 2} — 间隔不变性：$c^2\Delta t'^2 - \Delta x'^2 = c^2\Delta t^2 - \Delta x^2$
  \hfill\textit{(R1: kernel.lorentz\_invariance)}

  \textbf{Step 3} — 线性变换：$x' = Ax + Bt$，$t' = Cx + Dt$
  \hfill\texttt{[N]} 时空均匀 $\to$ 线性（无 $x^2$ 项）

  \textbf{Step 4} — 标准构型（S' 以速度 $v$ 相对 S 沿 $x$ 轴运动）

  \textbf{Step 5} — 代入间隔不变性，对照 $t^2$, $x^2$, $xt$ 系数：
  \begin{align*}
    t^2&: c^2D^2 - A^2v^2 = c^2\\
    x^2&: c^2C^2 - A^2 = -1\\
    xt&: 2c^2CD + 2A^2v = 0
  \end{align*}
  $\implies D = \gamma$，$A = \gamma$，$C = -\gamma v/c^2$，$B = -\gamma v$

  \textbf{Step 6} — $x' = \gamma(x-vt)$，$t' = \gamma(t-vx/c^2)$，$\gamma = 1/\sqrt{1-v^2/c^2}$
\end{derivationbox}

\texttt{reversible: true}，\texttt{[N]} $|v| < c$（$v=c$ 时 $\gamma$ 发散，$v>c$ 时 $\gamma$ 虚数——不物理）。

% ════════════════════════════════════════════════════════
\section{质能等价}
% ════════════════════════════════════════════════════════

\begin{lawbox}[质能等价]
  $E^2 = p^2c^2 + m^2c^4$；静止系：$E = mc^2$
  \quad\textbf{前提}：Lorentz 不变性 (R1)
\end{lawbox}

\noindent\textbf{变量}：四维动量 $p^\mu = (E/c, \mathbf{p})$ [M$\cdot$L$\cdot$T$^{-1}$],
静止质量 $m$ [M]（Lorentz 标量——与参考系无关的不变质量），
固有时 $\tau$ [T]（$d\tau^2 = ds^2/c^2$）。

\begin{derivationbox}[推导]
  \textbf{Step 1} — 四维动量：$p^\mu \equiv m\,dx^\mu/d\tau = (\gamma mc, \gamma m\mathbf{v}) = (E/c, \mathbf{p})$
  \hfill\textit{（$m$ 是 Lorentz 标量——粒子的内禀质量不随参考系改变）}

  \textbf{Step 2} — 时间分量：$E = \gamma mc^2$

  \textbf{Step 3} — Lorentz 不变量 (R1)：$p_\mu p^\mu = (E/c)^2 - \mathbf{p}^2 = m^2c^2$

  \textbf{Step 4} — $E^2 = \mathbf{p}^2c^2 + m^2c^4$。静止系（$\mathbf{p}=0$）：$E_{\text{rest}} = mc^2$
\end{derivationbox}

% ════════════════════════════════════════════════════════
\section{Einstein 场方程}
% ════════════════════════════════════════════════════════

\begin{lawbox}[Einstein 场方程]
  $G_{\mu\nu} + \Lambda g_{\mu\nu} = \dfrac{8\pi G}{c^4}\,T_{\mu\nu}$
  \quad\textbf{前提}：等效原理, 广义协变性 (R1), 最小作用量 (R2)
\end{lawbox}

\noindent\textbf{变量}：
Einstein 张量 $G_{\mu\nu} = R_{\mu\nu} - \frac{1}{2}Rg_{\mu\nu}$ [L$^{-2}$]，
Ricci 张量 $R_{\mu\nu}$ [L$^{-2}$]，
应力-能量张量 $T_{\mu\nu}$ [M$\cdot$L$^{-1}\cdot$T$^{-2}$]，
度规行列式 $g = \det(g_{\mu\nu}) < 0$ [1]。

\begin{derivationbox}[推导：Einstein-Hilbert 作用量]
  \textbf{Step 1} — Einstein-Hilbert 作用量：
  $S = \frac{1}{2\kappa}\int R\sqrt{-g}\,d^4x + S_{\text{matter}}$

  \textbf{注}：Lovelock 定理 (1971) 证明在 4D 中，EH 作用是唯一产生二阶运动方程的纯度规作用量。

  \textbf{Step 2} — 对 $g^{\mu\nu}$ 变分（R2：$\delta S = 0$）

  \textbf{Step 3} — 引力部分变分：$\delta(\int R\sqrt{-g}\,d^4x) = \int (R_{\mu\nu} - \frac{1}{2}Rg_{\mu\nu})\delta g^{\mu\nu}\sqrt{-g}\,d^4x$
  \hfill\textit{(Palatini 恒等式 + $\delta\sqrt{-g} = -\frac{1}{2}\sqrt{-g}g_{\mu\nu}\delta g^{\mu\nu}$)}

  \textbf{Step 4} — 定义 $G_{\mu\nu} \equiv R_{\mu\nu} - \frac{1}{2}Rg_{\mu\nu}$。$\nabla_\mu G^{\mu\nu} \equiv 0$ (Bianchi)

  \textbf{Step 5} — 物质部分：$T_{\mu\nu} \equiv -\frac{2}{\sqrt{-g}}\frac{\delta S_{\text{matter}}}{\delta g^{\mu\nu}}$

  \textbf{Step 6} — 场方程：$G_{\mu\nu} = \kappa T_{\mu\nu}$。Newton 极限（弱场低速）固定 $\kappa = 8\pi G/c^4$
\end{derivationbox}

\textbf{必要条件}：R1（等效原理 [N]——$m_i=m_g$ 使度规成为动力学场；广义协变性 [N]——作用量必须是标量），R2（$\delta S=0$ [N]），4 维时空 [N]（$G_{\mu\nu}$ 是唯一满足 $\nabla_\mu G^{\mu\nu}=0$ 的二阶张量）。

\textbf{充分条件}：R1+R2+4D $\Rightarrow$ Einstein 场方程（Lovelock 定理保证唯一性）。\texttt{reversible: true}

% ════════════════════════════════════════════════════════
\section{Newton 万有引力（弱场极限）}
% ════════════════════════════════════════════════════════

Newton 万有引力定律已在第 7 章作为 Einstein 场方程的弱场极限导出（第 189 页）。这里仅强调其推导链：

\begin{center}
  \textbf{R1（等效原理 + 广义协变性）+ R2 $\to$ Einstein 场方程 $\xrightarrow{\text{弱场} + \text{低速}}$ Newton 万有引力}
\end{center}

\texttt{reversible: false}——Newton 理论是 GR 的低能有效理论，丢失了时空弯曲的全部信息。
